\uuid{2UMD}
\titre{Noyau/Image/Injection/Surjection/Théorème du rang}
\auteur{hanani ; mzari}
\organisation{exo7}
\licence{CC BY-NC-SA 4.0}
\datecreate{2020-01-20}
\niveau{L1}
\module{Algèbre}
\chapitre{Application linéaire}
\sousChapitre{Définition}
\difficulte{2}
\theme{Noyau et image, Noyau, Image, Injection, Surjection, Théorème du rang}
\exo7qcmid{lille-2.applications-lineaires.015}

\contenu{
\begin{question}
\qclassification{107.01}
\qtags{motcle=Noyau/Image/Injection/Surjection/Théorème du rang}
On considère l' application linéaire :
$$\begin{array}{rccc}f:&\Rr^3&\to&\Rr^3\\
& (x,y,z)&\to &(x-z,y+z,x+y). \end{array}$$
Quelles sont les assertions vraies ?
\begin{answers}
\good{$\{(1,-1,1)\}$ est une base de $\ker f$}.
\bad{$f$ est injective}.
\good{$\{(1,0,1), (0,1,1)\}$ est une base de $\Im f$}.
\bad{$f$ est surjective}.
\end{answers}
\begin{explanations} $\ker f = \{(x,y,z) \in \Rr^3 \; ; \; (x-z,y+z,x+y)=(0,0,0)\} = \{(x,-x,x)  \; ; \; x\in \Rr\}$. Donc $ \{(1,-1,1)\}$ est une base de $\ker f$. Comme $\ker f \neq \{(0,0,0)\}$, $f$ n'est pas injective.
\vskip0mm
D'après le théorème du rang, $\dim \Im f = 2 $ et comme
$f(e_1)=(1,0,1)$ et $ f(e_2)=(0,1,1)$ ne sont pas colinéaires, ils forment une base de  $\Im f$. Comme $\dim \Im f = 2 $, $\Im f \neq \Rr^3 $, donc $f$ n'est pas surjective.
\end{explanations}
\end{question}
}
